\documentclass[a4paper]{article}
\usepackage{amsmath}
\usepackage{amssymb}
\usepackage{geometry}
\usepackage{hyperref}
\usepackage{booktabs}
\geometry{margin=1in}

\begin{document}

\title{\textit{Introduction to Econometrics}\\Assignment 2}
\author{Richard Harnisch, s5238366}
\date{\today}
\maketitle

% answer environment
\makeatletter
\newcommand{\answer}[1]{%
    \par\medskip
    \noindent
    \hspace*{-\@totalleftmargin}%
    \fbox{%
        \parbox{\dimexpr\textwidth-2\fboxsep-2\fboxrule}{#1}%
    }%
    \par\medskip
}
\makeatother

\noindent\textit{TeX code available under\\ \url{https://github.com/richardharnisch/intro-to-econometrics}}

\begin{enumerate}
    \item (3 points) Show that under assumptions 1.1--1.4 of Hayashi (2000), $\text{Cov}(\mathbf{b}, \mathbf{e} | \mathbf{X}) = \mathbf{0}$, where $\mathbf{e} = \mathbf{y} - \mathbf{Xb}$.

          \answer{
              By definition,
              \[
                  \text{Cov}(\mathbf{b}, \mathbf{e} | \mathbf{X}) = \mathbb{E}\{[\mathbf{b} - \mathbb{E}(\mathbf{b}|\mathbf{X})][\mathbf{e} - \mathbb{E}(\mathbf{e}|\mathbf{X})]' | \mathbf{X}\}.
              \]
              Under assumptions 1.1--1.4, the OLS estimator is $\mathbf{b} = \boldsymbol{\beta} + (\mathbf{X}'\mathbf{X})^{-1}\mathbf{X}'\boldsymbol{\varepsilon}$, and since $\mathbb{E}(\boldsymbol{\varepsilon}|\mathbf{X}) = \mathbf{0}$ (Assumption 1.2), we have $\mathbb{E}(\mathbf{b}|\mathbf{X}) = \boldsymbol{\beta}$, so
              \[
                  \mathbf{b} - \mathbb{E}(\mathbf{b}|\mathbf{X}) = (\mathbf{X}'\mathbf{X})^{-1}\mathbf{X}'\boldsymbol{\varepsilon}.
              \]
              For the residual vector, $\mathbf{e} = \mathbf{y} - \mathbf{Xb} = \mathbf{M}\mathbf{y} = \mathbf{M}(\mathbf{X}\boldsymbol{\beta} + \boldsymbol{\varepsilon}) = \mathbf{M}\boldsymbol{\varepsilon}$, where $\mathbf{M} = \mathbf{I} - \mathbf{X}(\mathbf{X}'\mathbf{X})^{-1}\mathbf{X}'$ is the annihilator matrix. Since $\mathbb{E}(\boldsymbol{\varepsilon}|\mathbf{X}) = \mathbf{0}$, we have $\mathbb{E}(\mathbf{e}|\mathbf{X}) = \mathbf{M}\mathbb{E}(\boldsymbol{\varepsilon}|\mathbf{X}) = \mathbf{0}$, so
              \[
                  \mathbf{e} - \mathbb{E}(\mathbf{e}|\mathbf{X}) = \mathbf{M}\boldsymbol{\varepsilon}.
              \]
              Substituting into the covariance formula:
              \begin{align*}
                  \text{Cov}(\mathbf{b}, \mathbf{e} | \mathbf{X}) & = \mathbb{E}\{[(\mathbf{X}'\mathbf{X})^{-1}\mathbf{X}'\boldsymbol{\varepsilon}][\mathbf{M}\boldsymbol{\varepsilon}]' | \mathbf{X}\}                                     \\
                                                                   & = (\mathbf{X}'\mathbf{X})^{-1}\mathbf{X}'\,\mathbb{E}(\boldsymbol{\varepsilon}\boldsymbol{\varepsilon}'|\mathbf{X})\,\mathbf{M}'.
              \end{align*}
              By Assumption 1.4, $\mathbb{E}(\boldsymbol{\varepsilon}\boldsymbol{\varepsilon}'|\mathbf{X}) = \sigma^2 \mathbf{I}$. Also, $\mathbf{M}' = \mathbf{M}$ (symmetric). Therefore,
              \begin{align*}
                  \text{Cov}(\mathbf{b}, \mathbf{e} | \mathbf{X}) & = \sigma^2 (\mathbf{X}'\mathbf{X})^{-1}\mathbf{X}'\mathbf{M}                                                                                                                \\
                                                                   & = \sigma^2 (\mathbf{X}'\mathbf{X})^{-1}\mathbf{X}'[\mathbf{I} - \mathbf{X}(\mathbf{X}'\mathbf{X})^{-1}\mathbf{X}']                                                           \\
                                                                   & = \sigma^2 [(\mathbf{X}'\mathbf{X})^{-1}\mathbf{X}' - (\mathbf{X}'\mathbf{X})^{-1}\mathbf{X}'\mathbf{X}(\mathbf{X}'\mathbf{X})^{-1}\mathbf{X}']                              \\
                                                                   & = \sigma^2 [(\mathbf{X}'\mathbf{X})^{-1}\mathbf{X}' - (\mathbf{X}'\mathbf{X})^{-1}\mathbf{X}']                                                                               \\
                                                                   & = \mathbf{0}.
              \end{align*}
          }

    \item (2 points) Consider the standard regression model $\mathbf{y} = \mathbf{X}\boldsymbol{\beta} + \boldsymbol{\varepsilon}$. Let $\hat{\boldsymbol{\beta}}$ be any estimator of $\boldsymbol{\beta}$. Let $\mathbf{u} = \hat{\boldsymbol{\beta}} - \mathbb{E}(\hat{\boldsymbol{\beta}}|\mathbf{X})$. Show that $\mathbb{E}[g(\mathbf{X})\mathbf{u}] = \mathbf{0}$ for any function $g(\mathbf{X})$.

          \answer{
              By the law of iterated expectations,
              \[
                  \mathbb{E}[g(\mathbf{X})\mathbf{u}] = \mathbb{E}\!\big[\mathbb{E}[g(\mathbf{X})\mathbf{u} | \mathbf{X}]\big].
              \]
              Since $g(\mathbf{X})$ is a function of $\mathbf{X}$ only, it can be pulled out of the inner conditional expectation:
              \[
                  \mathbb{E}[g(\mathbf{X})\mathbf{u} | \mathbf{X}] = g(\mathbf{X})\,\mathbb{E}[\mathbf{u} | \mathbf{X}].
              \]
              Now, $\mathbb{E}[\mathbf{u} | \mathbf{X}] = \mathbb{E}[\hat{\boldsymbol{\beta}} - \mathbb{E}(\hat{\boldsymbol{\beta}}|\mathbf{X}) | \mathbf{X}] = \mathbb{E}(\hat{\boldsymbol{\beta}}|\mathbf{X}) - \mathbb{E}(\hat{\boldsymbol{\beta}}|\mathbf{X}) = \mathbf{0}$.
              \\ \\
              Therefore, $\mathbb{E}[g(\mathbf{X})\mathbf{u}] = \mathbb{E}[g(\mathbf{X}) \cdot \mathbf{0}] = \mathbf{0}$.
          }

    \item Use the data in DISCRIM.DTA to answer this question.

          \begin{enumerate}
              \item (1 point) Estimate the model $\text{psoda}_i = \beta_1 + \beta_2 \text{prpblck}_i + \beta_3 \text{income}_i + \varepsilon_i$ by OLS.

                    \answer{
                        The OLS estimates are ($n = 401$, $R^2 = 0.064$):
                        \begin{center}
                            \begin{tabular}{lcccc}
                                \toprule
                                Variable & Coefficient & Std.\ Error & $t$-stat & $p$-value \\
                                \midrule
                                const    & 0.9563      & 0.019       & 50.354   & 0.000     \\
                                prpblck  & 0.1150      & 0.026       & 4.423    & 0.000     \\
                                income   & 1.603e--06  & 3.62e--07   & 4.430    & 0.000     \\
                                \bottomrule
                            \end{tabular}
                        \end{center}
                        \[
                            \widehat{\text{psoda}} = 0.9563 + 0.1150 \cdot \text{prpblck} + 0.000\,001\,603 \cdot \text{income}
                        \]
                    }

              \item (1 point) Give a meaningful economic interpretation of the constant term.

                    \answer{
                        The constant term $b_1 = 0.9563$ represents the estimated price of soda (in dollars) in a zip code area with zero percent African American population ($\text{prpblck} = 0$) and zero income ($\text{income} = 0$). Since zero income is not realistic, a more meaningful interpretation is that the constant anchors the regression line: it reflects the baseline soda price before accounting for the effects of racial composition and income. The estimated baseline price of approximately \$0.96 is broadly consistent with observed soda prices in the sample.
                    }

              \item (1 point) Consider an area with $\text{prpblck} = 0.4$. If the share of African Americans doubles, what is the estimated change in the price of soda?

                    \answer{
                        If $\text{prpblck}$ doubles from 0.4 to 0.8, the change is $\Delta\text{prpblck} = 0.4$. Since the model is linear, the estimated change in the price of soda is:
                        \[
                            \Delta\widehat{\text{psoda}} = b_2 \cdot \Delta\text{prpblck} = 0.1150 \times 0.4 = 0.0460.
                        \]
                        So the estimated price of soda increases by approximately \$0.046 (about 4.6 cents).
                    }

              \item (1 point) Estimate the model $\ln(\text{psoda}_i) = \beta_1 + \beta_2 \text{prpblck}_i + \beta_3 \ln(\text{income}_i) + \varepsilon_i$. What is the income elasticity of the price of soda?

                    \answer{
                        The OLS estimates are ($n = 401$, $R^2 = 0.068$):
                        \begin{center}
                            \begin{tabular}{lcccc}
                                \toprule
                                Variable       & Coefficient & Std.\ Error & $t$-stat & $p$-value \\
                                \midrule
                                const          & $-0.7938$   & 0.179       & $-4.424$ & 0.000     \\
                                prpblck        & 0.1216      & 0.026       & 4.722    & 0.000     \\
                                $\ln$(income)  & 0.0765      & 0.017       & 4.610    & 0.000     \\
                                \bottomrule
                            \end{tabular}
                        \end{center}
                        Since this is a log-log specification in income, the coefficient on $\ln(\text{income})$ is directly the income elasticity of the price of soda. The estimated income elasticity is $\hat{\beta}_3 = 0.0765$, meaning that a 1\% increase in income is associated with an approximately 0.077\% increase in the price of soda.
                    }

              \item (1 point) Estimate the model $\text{psoda}_i = \beta_1 + \beta_2 \text{prpblck}_i + \beta_3 \ln(\text{income}_i) + \varepsilon_i$. Provide a careful interpretation of $b_3$.

                    \answer{
                        The OLS estimates are ($n = 401$, $R^2 = 0.066$):
                        \begin{center}
                            \begin{tabular}{lcccc}
                                \toprule
                                Variable       & Coefficient & Std.\ Error & $t$-stat & $p$-value \\
                                \midrule
                                const          & 0.1855      & 0.188       & 0.987    & 0.324     \\
                                prpblck        & 0.1258      & 0.027       & 4.665    & 0.000     \\
                                $\ln$(income)  & 0.0788      & 0.017       & 4.533    & 0.000     \\
                                \bottomrule
                            \end{tabular}
                        \end{center}
                        This is a level-log specification. The coefficient $b_3 = 0.0788$ means that a 1\% increase in income is associated with an increase in the price of soda of approximately $0.0788 / 100 = 0.000788$ dollars, or about 0.08 cents, holding the proportion of African American population constant.
                    }
          \end{enumerate}
\end{enumerate}

\end{document}
